% Level 3, Week 1 answers -- Save Points: Git Basics
\answersweek{1}{Save Points: Git Basics}{Answers for workbook chapter 1 and homework 1.}

\begin{teachernote}[Placement quiz]
The placement quiz answers and scoring guide are in the assessments section at the end of this answer key.
\end{teachernote}

\section{Workbook exercises}
\begin{answer}[Exercise 1.1: The Git Quest]
There is no code to score: check the student's Terminal. In \texttt{Documents/git-quest} the commands
\texttt{git log --oneline} and \texttt{git status} should give something like:
\begin{outputcode}
3e8d1f2 (HEAD -> master) Make star bigger
9b0c7a4 Change badge color and hero name
7c41e0a Add hero badge drawing

On branch master
nothing to commit, working tree clean
\end{outputcode}
The branch may be called \texttt{main} instead of \texttt{master}; both are fine. An untracked \texttt{.DS\_Store} (a hidden Finder file) is harmless.
\textbf{What to look for:} 3 or more commits; messages that start with a verb and say what changed;
\texttt{git config user.name} shows a nickname and \texttt{git config user.email} an \texttt{@example.com} address;
the student can explain that \texttt{git restore} brought back the deleted \texttt{draw\_star} function.
The commit ids will always be different from these.
\end{answer}

\begin{answer}[Reflection: why two steps?]
\texttt{git add} lets you choose \emph{which} changes go into the next save point, so one commit can hold one idea (and a file with a secret can be left out). \texttt{git commit} then takes the snapshot and gives it a message. Any answer about choosing what to save, or checking before saving, is good.
\end{answer}

\subsection{Exercise 1.2: Commit Coach}
\pyfile{../exercises/week01/ex-commit-coach/solution.py}
\begin{answer}[Exercise 1.2: Commit Coach]
\textbf{What to look for:} the rules can be written in any order and with \texttt{and}/\texttt{or}, as long as all five checks pass.
In \texttt{find\_secrets}, the \texttt{break} stops a line being added twice when it has two secret words; a student who uses
\texttt{any()} or a flag variable is also correct.
\end{answer}

\subsection{Exercise 1.3: Log Reader}
\pyfile{../exercises/week01/ex-log-reader/solution.py}
\begin{answer}[Exercise 1.3: Log Reader]
\textbf{What to look for:} \texttt{split(" ", 1)} (split once) keeps messages with spaces in one piece. \texttt{oldest\_first} may use
\texttt{reverse()}, \texttt{reversed()} or slicing \texttt{[::-1]}. The last automatic check builds a real git repository in a temporary
folder; it is skipped on a Mac without git.
\end{answer}

\section{Homework}
\subsection{Core: Inventory Review}
\pyfile{../exercises/week01/hw-inventory-review/solution.py}
\begin{answer}[Core: Inventory Review]
Expected output includes \texttt{Gems: 3}, \texttt{Most common: gem} and \texttt{After using all potions: ['gem', 'key', 'gem', 'map', 'gem']}.
Common slip: removing items from the list while looping over it (it skips items) instead of building a new list.
The git step should add one more commit (\enquote{Add inventory review} or similar), so \texttt{git-quest} now has at least 4 commits, and the file in it is called \texttt{inventory\_review.py}.
\end{answer}

\begin{answer}[Challenge: Message makeover]
Any message that starts with a verb, says what changed and has 50 or fewer characters. Examples:
\enquote{Make hero jump higher}, \enquote{Add sound when you score}, \enquote{Stop score going below zero}.
All three should get \enquote{Great message!} from the Commit Coach.
\end{answer}

\begin{answer}[Creative: A badge with a history]
\textbf{Good looks like:} at least 3 new commits, one change per commit, a timeline drawing whose circles match \texttt{git log --oneline}
(oldest on the left). The hero name is invented. Praise small, focused commits more than one big one.
\end{answer}
